\documentclass[11pt]{amsart}

\usepackage[T1]{fontenc}
\usepackage{lmodern}
\usepackage{microtype}
\usepackage{mathtools,amssymb,amsthm}
\usepackage[hidelinks]{hyperref}

\emergencystretch=1.5em

\hypersetup{
  pdftitle={The Clow-Zaguia cop bound is tight at independence number four: f(4)=3}
}

\newtheorem{theorem}{Theorem}
\newtheorem{lemma}[theorem]{Lemma}
\theoremstyle{definition}
\newtheorem*{problemCZ}{Problem 5.2 (Clow and Zaguia)}

\DeclareMathOperator{\cop}{c}
\newcommand{\al}{\alpha}
\newcommand{\theonum}{\theta}
\newcommand{\Gwit}{G}

\title[The cop bound is tight at $\al=4$]
{The Clow--Zaguia Cop Bound is Tight at Independence Number Four: $f(4)=3$}

\date{}

\subjclass[2020]{05C57, 05C17, 05C69, 91A43}
\keywords{cops and robbers, cop number, perfect graph, independence number,
clique cover, rook's graph, tensor product, Latin square graph}

\begin{document}

\begin{abstract}
Clow and Zaguia proved that every connected perfect graph $G$ with
$\al(G)\ge 4$ satisfies $\cop(G)<\al(G)$, and asked for the least function
$f$ with $\cop(G)\le f(\al(G))$ on connected perfect graphs. We settle the
value $f(4)=3$. The witness is the complement of the $4\times4$ rook's graph,
equivalently the complement of $L(K_{4,4})$, equivalently the tensor product
$K_4\times K_4$: a $9$-regular graph on $16$ vertices that is perfect with
$\al=\theonum=4$ and cop number exactly $3$. Both bounds are proved by hand on
the printed object; the lower bound $\cop\ge3$ is an explicit robber strategy
against two cops, given as a six-case invariant. The graph and its cop number
are already in print, in work of Ahirwar, Bonato, Gittins, Huang, Marbach and
Zaidman on Latin square graphs; perfection together with $\al=4$ identifies it
as a witness for $f(4)$.
\end{abstract}

\maketitle

\section{The problem}

We use the game of Cops and Robbers in the convention of Clow and
Zaguia~\cite{CZ}: the cops occupy vertices first, the robber then chooses a
vertex, and thereafter the two sides move alternately with the cops moving
first, each token either staying where it is or moving to an adjacent vertex.
The cops win if some cop ever occupies the robber's vertex. The cop number
$\cop(G)$ is the least number of cops with a winning strategy. As
usual $\al(G)$ is the independence number and $\theonum(G)$ the clique cover
number. \emph{Perfect} is used in the form Clow and Zaguia state
\cite[\S1]{CZ}, which by the Strong Perfect Graph Theorem of Chudnovsky,
Robertson, Seymour and Thomas~\cite{SPGT} is equivalent to the usual one: $G$
contains no induced odd hole $C_{2k+1}$ and no induced odd antihole,
$k\ge2$.

The following is Problem 5.2 of~\cite{CZ}, quoted from the arXiv version of
that paper.

\begin{problemCZ}
Given an integer $k$, what is the least integer $f(k)$ such that every
connected perfect graph $G$ with $\al(G)=k$, has $\cop(G)\le f(k)$?
\end{problemCZ}

A witness must satisfy all three clauses, and all three are checked below: $G$
connected, $G$ perfect, and $\al(G)$ equal to $k$ \emph{exactly}, not merely at
least $k$.

The upper half of the case $k=4$ is a corollary of the same paper, which
proves:

\begin{quote}
If $G$ is a connected perfect graph with $\al(G)\ge4$, then
$\cop(G)<\al(G)$.
\end{quote}

Hence $f(k)\le k-1$ for every $k\ge4$, and in particular $f(4)\le3$. What the
$k=4$ instance of Problem 5.2 therefore needs is a single connected perfect
graph with $\al=4$ and cop number $3$. We exhibit one, so the bound of
\cite{CZ} is tight at $k=4$, and $f(4)=3$.

\section{The witness}\label{sec:witness}

Let $\Gwit$ be the graph whose vertices are the $16$ cells $(i,j)$ of a
$4\times4$ grid, $i,j\in\{1,2,3,4\}$, with
\begin{equation}\label{eq:rule}
  (i,j)\sim(k,l)
  \quad\Longleftrightarrow\quad
  i\neq k \ \text{ and }\ j\neq l .
\end{equation}
Say two distinct cells \emph{share a line} if they agree in their first or in
their second coordinate. Then \eqref{eq:rule} says exactly that two distinct
cells are non-adjacent if and only if they share a line, which is the only
fact about $\Gwit$ used below. Equivalently $\Gwit$ is the complement of the
$4\times4$ rook's graph $K_4\mathbin{\square}K_4=L(K_{4,4})$, and also the
tensor (categorical) product $K_4\times K_4$.

The non-neighbours of $(i,j)$ are the other three cells of row $i$ and the
other three cells of column $j$, so $\Gwit$ is $9$-regular with
$|E(\Gwit)|=\tfrac{16\cdot9}{2}=72=\binom{16}{2}-48$, the $48$ non-edges being
the $4$ rows and $4$ columns at $\binom42=6$ pairs each. Indexing the cell
$(i,j)$ as vertex $4(i-1)+(j-1)$ and listing cells in lexicographic order
$11,12,13,14,21,\dots,44$, the \texttt{graph6} string of $\Gwit$ is
\begin{center}
\texttt{O?]uf@vmuy\textbackslash o\textbackslash vmzZmf\textbackslash o}
\end{center}
The proofs below check two further objects, printed here.
\begin{itemize}
\item The four \emph{broken diagonals}
  $Q_t=\{(i,\,1+((i-1+t)\bmod 4)):1\le i\le4\}$ for $t=0,1,2,3$, that is
  \[
    \begin{array}{l}
      Q_0=\{(1,1),(2,2),(3,3),(4,4)\},\quad
      Q_1=\{(1,2),(2,3),(3,4),(4,1)\},\\[2pt]
      Q_2=\{(1,3),(2,4),(3,1),(4,2)\},\quad
      Q_3=\{(1,4),(2,1),(3,2),(4,3)\}.
    \end{array}
  \]
\item The three-cell set $D=\{(1,1),(2,2),(3,3)\}$.
\end{itemize}

\begin{lemma}\label{lem:basic}
$\Gwit$ is connected, perfect, and $\al(\Gwit)=\theonum(\Gwit)=4$.
\end{lemma}

\begin{proof}
\emph{Connected.} Two cells sharing a line have a common
neighbour, namely any cell off both of their rows and both of their columns;
two cells occupy at most $2$ rows and at most $2$ columns, so at least
$4-2=2$ rows and at least $2$ columns remain and such a cell exists.

\emph{Perfect.} The rook's graph $K_4\mathbin{\square}K_4$ is the line graph
$L(K_{4,4})$ of a bipartite graph --- the cell $(i,j)$ corresponds to the edge
of $K_{4,4}$ joining the $i$th vertex of one side to the $j$th vertex of the
other, and two cells share a line precisely when the two edges share an
endpoint. Line graphs of bipartite graphs are perfect, by K\"onig's edge
colouring theorem; and the complement of a perfect graph is perfect, by the
Perfect Graph Theorem of Lov\'asz~\cite{Lovasz72}, which \cite[\S1]{CZ} also
quotes. Hence $\Gwit=\overline{L(K_{4,4})}$ is perfect.

\emph{$\al=4$.} An independent set $S$ of $\Gwit$ is a set of cells that
pairwise share a line. Suppose $a,b\in S$ lie in the same row $i$ with
distinct columns $j_1\neq j_2$. Any $c\in S$ off row $i$ must share a line
with both, hence must lie in column $j_1$ and in column $j_2$, which is
impossible. So $S$ lies inside a single row; symmetrically, if two elements of
$S$ share a column then $S$ lies inside a single column. Either way
$|S|\le4$. A full row is independent, so $\al(\Gwit)=4$.

\emph{$\theonum=4$.} Each $Q_t$ has all four rows distinct and all four
columns distinct, so by \eqref{eq:rule} it is a clique, indeed a $K_4$; and
$Q_0,Q_1,Q_2,Q_3$ partition the $16$ cells. Hence $\theonum\le4$. Also
$\theonum\ge\al=4$, so $\theonum=4$.
\end{proof}

\begin{lemma}\label{lem:upper}
$\cop(\Gwit)\le3$.
\end{lemma}

\begin{proof}
$D=\{(1,1),(2,2),(3,3)\}$ is a dominating set: a cell $(i,j)$ undominated by
$D$ would have to satisfy $i=1$ or $j=1$, and $i=2$ or $j=2$, and $i=3$ or
$j=3$; since $i$ equals at most one of $1,2,3$, the value $j$ would have to
equal two distinct elements of $\{1,2,3\}$. Three cops placed on a dominating
set capture on their first move, so $\cop(\Gwit)\le\gamma(\Gwit)\le3$.
\end{proof}

The lower bound is the substance. It is a strategy for the robber against two
cops, kept honest by a single invariant.

\begin{lemma}\label{lem:lower}
Let $m\ge4$ and let $\Gwit_m=K_m\times K_m$, that is, the graph on the $m^2$
cells of an $m\times m$ grid with the adjacency rule \eqref{eq:rule}. Then two
cops do not suffice to catch the robber on $\Gwit_m$; hence
$\cop(\Gwit_m)\ge3$. In particular $\cop(\Gwit)\ge3$.
\end{lemma}

\begin{proof}
Call a cell $r$ \emph{safe} for a pair of cop positions $(u_1,u_2)$ if $r$
shares a line with $u_1$ and with $u_2$. By \eqref{eq:rule} this says exactly
that $r\notin\{u_1,u_2\}$ and $r$ is adjacent to neither cop. We prove three
statements; together they say that the robber can occupy a safe cell after
every one of his moves, and therefore is never caught.

\smallskip
\noindent\emph{(C) Opening.} For every pair $(u_1,u_2)=((a,b),(c,d))$ of cop
starting cells, equal cells allowed, a safe cell exists, and the following rule
names one. If $a=c$, take $r=(a,e)$ with $e\notin\{b,d\}$, which exists since
$m\ge3>|\{b,d\}|$; then $r$ lies in row $a$ with both cops and equals neither.
Otherwise, if $b=d$, take $r=(e,b)$ with $e\notin\{a,c\}$, symmetrically.
Otherwise $a\neq c$ and $b\neq d$; take $r=(a,d)$, which shares row $a$ with
$u_1$ and column $d$ with $u_2$, and differs from $u_1$ in the second
coordinate and from $u_2$ in the first.

\smallskip
\noindent\emph{(A) A cop move out of a safe position cannot capture.} Let $r$
be safe for $(u_1,u_2)$ and let the cops move to $(u_1',u_2')$ with
$u_i'\in N[u_i]$. If $u_i'=r$ then either $u_i'=u_i=r$, contradicting
$r\notin\{u_1,u_2\}$, or $u_i\sim u_i'=r$, contradicting that $r$ is
non-adjacent to $u_i$.

\smallskip
\noindent\emph{(B) Safety can be restored.} Let $r=(x,y)$ be safe for the cop
positions before their move, and let the cops now stand at $u_1'=(p,q)$ and
$u_2'=(s,t)$; by (A), $r\notin\{u_1',u_2'\}$. We exhibit
$r'\in N[r]$ safe for $(u_1',u_2')$; note that $N[r]$ consists of $r$ together
with all cells $(u,v)$ with $u\neq x$ and $v\neq y$. Say a cop
\emph{threatens} $r$ if it is adjacent to $r$, i.e.\ $u_1'$ threatens $r$
exactly when $p\neq x$ and $q\neq y$.
\begin{description}
\item[Case 0, neither cop threatens $r$.] Then $r$ itself is safe for the new
  positions; the robber passes, $r'=r$.
\item[Case 1a, both threaten and $p=s$.] Take $r'=(p,v)$ with
  $v\notin\{y,q,t\}$, possible since $m\ge4>|\{y,q,t\}|$. Then $r'$ is in row
  $p$, so it shares a line with both cops; $p\neq x$ and $v\neq y$, so
  $r'\in N[r]$; and $r'$ is neither $(p,q)$ nor $(s,t)=(p,t)$.
\item[Case 1b, both threaten, $p\neq s$ and $q\neq t$.] Take $r'=(p,t)$: it
  shares row $p$ with $u_1'$ and column $t$ with $u_2'$; $p\neq x$ and
  $t\neq y$; and $r'\neq(p,q)$ as $t\neq q$, $r'\neq(s,t)$ as $p\neq s$.
\item[Case 1c, both threaten, $p\neq s$ and $q=t$.] Take $r'=(u,q)$ with
  $u\notin\{x,p,s\}$, possible since $m\ge4>|\{x,p,s\}|$. Then $r'$ is in
  column $q=t$, so it shares a line with both cops; $u\neq x$ and $q\neq y$;
  and $r'$ is neither $(p,q)$ nor $(s,q)$.
\item[Cases 2a and 2b, exactly one cop threatens $r$.] Relabel so that
  $u_1'=(p,q)$ threatens $r$ and $u_2'=(s,t)$ does not, so $s=x$ or $t=y$; not
  both, since $(s,t)=(x,y)=r$ is excluded by (A).
  \begin{description}
  \item[Case 2a, $s=x$ and $t\neq y$.] Any $r'=(u,v)\in N[r]\setminus\{r\}$
    has $u\neq x=s$, so sharing a line with $u_2'$ forces $v=t$; and $t\neq y$
    is given. If $t\neq q$ take $r'=(p,t)$, which shares row $p$ with $u_1'$
    and column $t$ with $u_2'$, and differs from $(p,q)$ and from $(x,t)$. If
    $t=q$ take $r'=(u,t)$ with $u\notin\{x,p\}$, possible since $m\ge3$; then
    $r'$ shares column $t=q$ with $u_1'$ and column $t$ with $u_2'$, and
    equals neither.
  \item[Case 2b, $t=y$ and $s\neq x$.] Symmetrically any
    $r'=(u,v)\in N[r]\setminus\{r\}$ has $v\neq y=t$, so $u=s$. If $s\neq p$
    take $r'=(s,q)$: it shares column $q$ with $u_1'$ and row $s$ with
    $u_2'$, and equals neither. If $s=p$ take $r'=(s,v)$ with
    $v\notin\{y,q\}$, possible since $m\ge3$; then $r'$ is in row $s=p$ with
    both cops and equals neither.
  \end{description}
\end{description}
No case imposes any condition on \emph{which} move the cops made, so the
argument covers every cop strategy. By (C) the robber can start safe; by (A)
no cop move ever captures him; by (B) he restores safety after each cop move.
Hence he survives forever and $\cop(\Gwit_m)\ge3$. Case 0 uses the robber's
option to pass, which is part of the convention of~\cite{CZ}.
\end{proof}

\begin{theorem}\label{thm:main}
$\Gwit$ is a connected perfect graph with $\al(\Gwit)=4$ and
$\cop(\Gwit)=3$. Consequently $f(4)=3$: the bound $\cop(G)<\al(G)$ of
\cite{CZ} for connected perfect graphs with $\al\ge4$ is tight at $\al=4$.
\end{theorem}

\begin{proof}
Lemma~\ref{lem:basic} gives connectivity, perfection and $\al=4$;
Lemmas~\ref{lem:upper} and~\ref{lem:lower} give $3\le\cop(\Gwit)\le3$. So
$f(4)\ge3$. The corollary of \cite{CZ} quoted in \S1 gives $f(4)\le3$.
\end{proof}

Only the instance $k=4$ of Problem 5.2 is settled here; no value $f(k)$ with
$k\neq4$ is determined.

\section{Attribution}\label{sec:prior}

The graph $\Gwit$ and the value $\cop(\Gwit)=3$ are already in print. Ahirwar,
Bonato, Gittins, Huang, Marbach and Zaidman~\cite{Latin} study Cops and
Robbers on \emph{Latin square graphs}: for an $m\times m$ Latin square $L$,
the graph $G(L)$ has the $m^2$ cells as vertices, two distinct cells being
adjacent when they share a row, share a column, or carry the same symbol.
They record that ``by directly checking, the cop number of a latin square of
order 1 or 2 is 1, order 3 is 2, and order 4 is 3''.

Our $\Gwit$ is such a graph. Let $L$ be the Cayley table of the Klein
four-group, $L(i,j)=i\oplus j$ on $\{0,1,2,3\}$ with $\oplus$ bitwise
exclusive or, and index the cells of $G(L)$ as $4i+j$, $0\le i,j\le3$. Then
the permutation
\[
  \varphi=(0,\,5,\,10,\,15,\,11,\,14,\,1,\,4,\,13,\,8,\,7,\,2,\,6,\,3,\,12,\,9),
\]
read as $\varphi(k)$ for $k=0,\dots,15$ with vertex $k$ of $\Gwit$ being the
cell $(1+\lfloor k/4\rfloor,\,1+(k\bmod4))$, is an isomorphism
$G(L)\to\Gwit$; a reader can check the $\binom{16}{2}=120$ pairs directly. So
the object of Theorem~\ref{thm:main} is the order-$4$ Latin square graph of
\cite{Latin}, and the integer $3$ is theirs. What perfection and
$\al(\Gwit)=4$ add is that this graph is a witness for $f(4)$.

\section{Verification}

Every claim above is checkable by hand on the objects printed in
\S\ref{sec:witness} and \S\ref{sec:prior}. A program accompanies the note,
\texttt{verify.py}, with its recorded output in \texttt{verify.output.txt}; it
uses exact integer and bitmask arithmetic and no third-party package. It reads
the \texttt{graph6} string, the adjacency rule, the broken diagonals, the
dominating set, the six-case strategy and the permutation $\varphi$ as printed
here; it recomputes $\al$, $\theonum$ and the domination number by exact
search; it certifies perfection twice, once by enumerating all $32{,}192$ odd
vertex subsets of size at least $5$ for an induced odd hole or antihole and
once by checking that the complement is $L(K_{4,4})$; it replays
Lemma~\ref{lem:lower} over all $576$ safe states and all $57{,}600$ cop moves
out of them, in the case-by-case form printed above, and independently decides
the two-cop game by backward induction over all $16^3\cdot2$ states; and it
verifies $\varphi$. Its recorded run reports
\begin{center}
\texttt{VERDICT: ALL 37 CHECKS PASS}
\end{center}
No three-cop solver was run: $\cop(\Gwit)\le3$ rests on the printed dominating
set, and independently on the corollary of \cite{CZ} quoted in \S1.

\begin{thebibliography}{9}

\bibitem{Latin}
S.~Ahirwar, A.~Bonato, L.~Gittins, A.~Huang, T.~G. Marbach, and T.~Zaidman,
\emph{Pursuit-evasion games on latin square graphs},
J. Comb. \textbf{14} (2023), no.~4, 461--483.
\href{https://doi.org/10.4310/JOC.2023.v14.n4.a4}{doi:10.4310/JOC.2023.v14.n4.a4};
\href{https://arxiv.org/abs/2109.14669}{arXiv:2109.14669}.
The sentence quoted in \S\ref{sec:prior} is from the arXiv e-print.

\bibitem{SPGT}
M.~Chudnovsky, N.~Robertson, P.~Seymour, and R.~Thomas,
\emph{The strong perfect graph theorem},
Ann. of Math. (2) \textbf{164} (2006), no.~1, 51--229.
\href{https://doi.org/10.4007/annals.2006.164.51}{doi:10.4007/annals.2006.164.51}.

\bibitem{CZ}
A.~Clow and I.~Zaguia,
\emph{Cops and Robbers, clique covers, and induced cycles},
Discrete Math. \textbf{349} (2026), 115252.
\href{https://doi.org/10.1016/j.disc.2026.115252}{doi:10.1016/j.disc.2026.115252};
\href{https://arxiv.org/abs/2507.14321}{arXiv:2507.14321v1}.
Problem 5.2 and the corollary quoted in \S1 are reproduced in full above from
the arXiv source, version v1, which is the text we read; the numbering of the
journal text may differ.

\bibitem{Lovasz72}
L.~Lov\'asz,
\emph{Normal hypergraphs and the perfect graph conjecture},
Discrete Math. \textbf{2} (1972), no.~3, 253--267.
\href{https://doi.org/10.1016/0012-365X(72)90006-4}{doi:10.1016/0012-365X(72)90006-4}.

\end{thebibliography}

\end{document}
